\section*{Chapter \ref{ch:b2:hermitian} --- Hermitian Forms}

\begin{solution}{exo:b2:hermitian:1}
$\langle x, y\rangle = \conj{1}\cdot\iu + \conj{\iu}\cdot 1 = \iu -
\iu = 0$: orthogonal. $\norm x = \norm y = \sqrt{1 + 1} = \sqrt2$.
Orthonormal basis: $\bigl(\frac{1}{\sqrt2}(1, \iu),\;
\frac{1}{\sqrt2}(\iu, 1)\bigr)$ --- the normalized pair itself.
\end{solution}

\begin{solution}{exo:b2:hermitian:2}
First: equals its conjugate transpose (real diagonal, $\conj{\iu} =
-\iu$ swapped): Hermitian (hence normal); not unitary ($A^\dagger A
\neq I$: columns not unit).

Second: $A^\dagger A = \frac12\begin{pmatrix} 1 & -\iu\\ -\iu &
1\end{pmatrix}\begin{pmatrix} 1 & \iu\\ \iu & 1\end{pmatrix} =
\frac12\begin{pmatrix} 2 & 0\\ 0 & 2\end{pmatrix} = I$: unitary
(hence normal); not Hermitian.

Third: $A^\dagger A = E_{22} \neq E_{11} = AA^\dagger$: not normal
(so neither Hermitian nor unitary) --- the standard nilpotent
counterexample.
\end{solution}

\begin{solution}{exo:b2:hermitian:3}
Uniqueness: $A = H + \iu K$ with $H^\dagger = H$, $K^\dagger = K$
forces $A^\dagger = H - \iu K$, so $H = \frac{A + A^\dagger}{2}$,
$K = \frac{A - A^\dagger}{2\iu}$; these formulas are Hermitian
(check: $\bigl(\frac{A - A^\dagger}{2\iu}\bigr)^\dagger = \frac{
A^\dagger - A}{-2\iu} = K$) and reconstruct $A$: existence.

Normality: $A^\dagger A - AA^\dagger = (H - \iu K)(H + \iu K) - (H
+ \iu K)(H - \iu K) = 2\iu(HK - KH)$: it vanishes iff $HK = KH$.
\end{solution}

\begin{solution}{exo:b2:hermitian:4}
$\chi_A = (X-2)^2 - 1$: eigenvalues $3$ and $1$. Eigenvectors, by
direct computation:
\[
A\begin{pmatrix}1\\ -\iu\end{pmatrix}
= \begin{pmatrix} 2 + \iu(-\iu)\\ -\iu + 2(-\iu)\end{pmatrix}
= \begin{pmatrix} 3\\ -3\iu \end{pmatrix}
= 3\begin{pmatrix}1\\ -\iu\end{pmatrix},
\qquad
A\begin{pmatrix}1\\ \iu\end{pmatrix}
= \begin{pmatrix} 2 + \iu\cdot\iu\\ -\iu + 2\iu\end{pmatrix}
= \begin{pmatrix}1\\ \iu\end{pmatrix}.
\]
Orthonormal eigenbasis: $u_1 = \frac{1}{\sqrt2}(1, -\iu)$
(eigenvalue $3$), $u_2 = \frac{1}{\sqrt2}(1, \iu)$ (eigenvalue
$1$); orthogonality as in \cref{exo:b2:hermitian:1}. Powers, via
the spectral projections $A^k = 3^k u_1u_1^\dagger + 1^k\,
u_2u_2^\dagger$:
\[
A^k = U\begin{pmatrix} 3^k & 0\\ 0 & 1\end{pmatrix}U^\dagger
= \frac{3^k}{2}\begin{pmatrix} 1 & \iu\\ -\iu & 1\end{pmatrix}
+ \frac{1}{2}\begin{pmatrix} 1 & -\iu\\ \iu & 1\end{pmatrix}
= \frac12\begin{pmatrix}
3^k + 1 & (3^k - 1)\iu\\
-(3^k-1)\iu & 3^k + 1
\end{pmatrix}.
\]
(Check $k = 1$: recovers $A$.)
\end{solution}

\begin{solution}{exo:b2:hermitian:5}
Compact: closed (preimage of $I$ under the continuous $U \mapsto
U^\dagger U$) and bounded (columns are unit vectors: entries of
modulus $\leq 1$) in $\mathcal{M}_n(\C) \simeq \R^{2n^2}$.

Eigenvalues: $\operatorname{diag}(\lambda, 1, \dots, 1)$ is unitary
for any $\abs\lambda = 1$. Determinant: $\abs{\det U}^2 = \det
U^\dagger \det U = \det(U^\dagger U) = 1$ (using $\det A^\dagger =
\conj{\det A}$): $\det U$ lies on the unit circle.
\end{solution}

\begin{solution}{exo:b2:hermitian:6}
By the spectral theorem, work in an orthonormal eigenbasis of $H$:
everything reduces to scalars $\lambda \in \R$ (the eigenvalues).
$H + \iu I$ has eigenvalues $\lambda + \iu \neq 0$: invertible. $U$
has eigenvalues $\mu = \frac{\lambda - \iu}{\lambda + \iu}$, of
modulus $1$ ($\abs{\lambda - \iu} = \abs{\lambda + \iu}$ for real
$\lambda$): $U^\dagger U = I$ holds since $U$ is unitarily
diagonalizable with unimodular eigenvalues (it is diagonal in the
chosen orthonormal basis). And $\mu = 1$ would force $-\iu = \iu$:
impossible, so $1 \notin \operatorname{Sp} U$. (The Cayley
transform maps Hermitian to unitary-minus-a-point --- the matrix
version of the map from $\R$ to the circle.)
\end{solution}

\begin{solution}{exo:b2:hermitian:7}
For an eigenpair $Ax = \lambda x$ ($x \neq 0$): $\lambda\norm x^2 =
\langle x, Ax\rangle > 0$, so $\lambda > 0$ (already real,
\cref{prop:b2:hermitian:eigenvalues}). Square root: in a spectral
basis, $B = U\operatorname{diag}(\sqrt{\lambda_i})U^\dagger$:
Hermitian, positive definite, $B^2 = A$. Determinant: product of
the positive eigenvalues.
\end{solution}

\begin{solution}{exo:b2:hermitian:8}
\begin{enumerate}
  \item $\norm{u(x)}^2 = \langle u(x), u(x)\rangle = \langle x,
        u^*u(x)\rangle = \langle x, uu^*(x)\rangle =
        \norm{u^*(x)}^2$. Applying this to the normal $u -
        \lambda\,\mathrm{id}$ (its adjoint is $u^* -
        \conj\lambda$, and normality is inherited): $\norm{(u -
        \lambda)x} = \norm{(u^* - \conj\lambda)x}$, so the kernels
        agree.
  \item For eigenvectors $u(x) = \lambda x$, $u(y) = \mu y$
        ($\lambda \neq \mu$): using (1), $u^*(x) = \conj\lambda x$;
        then
        \[
        \lambda\langle y, x\rangle = \langle y, u(x)\rangle
        = \langle u^*(y), x\rangle = \langle \conj\mu\, y, x\rangle
        = \mu \langle y, x\rangle ,
        \]
        so $\langle y, x\rangle = 0$. Stability of $E_\lambda^\perp$:
        for $x \perp E_\lambda$ and $z \in E_\lambda$, $\langle z,
        u(x)\rangle = \langle u^*(z), x\rangle =
        \langle\conj\lambda z, x\rangle = 0$.
  \item Induction on dimension: over $\C$, $u$ has an eigenvector
        $e_1$ (normalize); its orthogonal complement is stable
        under $u$ (by (2)) \emph{and} under $u^*$ (same argument
        with roles swapped), so the restriction is normal: induct
        and concatenate orthonormal eigenbases. Conversely, a
        unitarily diagonalizable $u = UDU^\dagger$ satisfies $u^*u
        = U\conj D D U^\dagger = U D\conj D U^\dagger = uu^*$:
        normal.
\end{enumerate}
\end{solution}
